\begin{figure}[H]
\centering
\resizebox{0.96\linewidth}{!}{
\begin{tikzpicture}[
    >=Latex,
    font=\small,
    every path/.style={line width=0.9pt},
    mat/.style={draw, rounded corners=2pt, minimum width=16mm, minimum height=32mm, inner sep=2pt, fill=gray!10},
    block/.style={draw, rounded corners=3pt, minimum width=22mm, minimum height=10mm, align=center, fill=orange!15},
    opblock/.style={draw, rounded corners=3pt, minimum width=24mm, minimum height=10mm, align=center, fill=blue!12},
    tiny/.style={font=\scriptsize},
    lbl/.style={font=\small}
]

% Left: input embeddings H^(l)
\node[mat] (Hin) at (0,0) {};
\draw (-0.8,0.8) -- (0.8,0.8);
\draw (-0.8,0.4) -- (0.8,0.4);
\draw (-0.8,0.0) -- (0.8,0.0);
\draw (-0.8,-0.4) -- (0.8,-0.4);
\draw (-0.8,-0.8) -- (0.8,-0.8);
\node[lbl] at (0,-2.1) {$\TRFH^{(l)} \in \mathbb{R}^{n\times d}$};

% Q K V projections
\node[opblock] (qproj) at (3.2,1.6) {$\TRFQ = \TRFH^{(l)}\TRFWQ$};
\node[opblock] (kproj) at (3.2,0.0) {$\TRFK = \TRFH^{(l)}\TRFWK$};
\node[opblock] (vproj) at (3.2,-1.6) {$\TRFV = \TRFH^{(l)}\TRFWV$};

\draw[->] (Hin.east) -- ++(0.7,0) |- (qproj.west);
\draw[->] (Hin.east) -- ++(0.7,0) -- (kproj.west);
\draw[->] (Hin.east) -- ++(0.7,0) |- (vproj.west);

% attention formula block
\node[block, minimum width=35mm, minimum height=16mm] (attn) at (7.4,0){$\ATTN(\TRFH^{(l)})$};

\draw[->] (qproj.east) -- ++(0.8,0) |- (attn.west);
\draw[->] (kproj.east) -- (attn.west);
\draw[->] (vproj.east) -- ++(0.8,0) |- (attn.west);



% First residual + norm
\node[block] (addnorm1) at (11.2,0) {Add \& Norm};

%\draw[->] (attn.east) -- (mha.west);
\draw[->] ($(attn.east)+(0,-0.2)$) -- ($(addnorm1.west)+(0,-0.2)$);

% residual from H^(l)
\draw[->] (Hin.north) |- ++(0,1.2) -| ($(addnorm1.west)+(-0.4,0.2)$) -- ($(addnorm1.west)+(0,0.2)$);

\node[lbl] at (13.2,-2.1) {$\tilde \TRFH^{(l)}$};

% FFN block
\node[block] (ffn) at (14.4,0) {FFN};

% Second residual + norm
\node[block] (addnorm2) at (17.6,0) {Add \& Norm};

\draw[->] ($(addnorm1.east)+(0,-0.2)$) -- ($(ffn.west)+(0,-0.2)$);
\draw[->] 
    ($(addnorm1.east)+(0,0.2)$) 
    -- ++(0.4,0)
    -- ++(0,2.6)
    -| ($(addnorm2.west)+(-0.4,0.2)$) 
    -- ($(addnorm2.west)+(0,0.2)$);

\draw[->] ($(ffn.east)+(0,-0.2)$) -- ($(addnorm2.west)+(0,-0.2)$);

% Output H^(l+1)
\node[mat] (Hout) at (21.0,0) {};
\draw (20.2,0.8) -- (21.8,0.8);
\draw (20.2,0.4) -- (21.8,0.4);
\draw (20.2,0.0) -- (21.8,0.0);
\draw (20.2,-0.4) -- (21.8,-0.4);
\draw (20.2,-0.8) -- (21.8,-0.8);

\draw[->] (addnorm2.east) -- (Hout.west);

\node[lbl] at (21.0,-2.1) {$\TRFH^{(l+1)} \in \mathbb{R}^{n\times d}$};

\end{tikzpicture}}
\caption{Generic transformer layer with a single self-attention head. Starting from input embeddings $\TRFH^{(l)}$, learned projections produce queries, keys, and values. Self-attention aggregates information across the input elements, after which residual connections, normalization, and a position-wise feed-forward network produce the updated embeddings $\TRFH^{(l+1)}$.}
\label{fig:background_transformer_layer}
\end{figure}